% Contexte, objectifs, problématique, contraintes (rapport de cadrage, A et C).

\chapter{Introduction}
\label{chap:introduction}

\section{Contexte}

Le projet est né d'un dépôt existant : un modèle de Poisson qui fonctionnait, une collecte de données
depuis l'API commerciale API-FOOTBALL, une première application. Le rapport de cadrage du 24 septembre
2026 l'a analysé et a recommandé une \emph{refonte ciblée} du cœur de données plutôt qu'une réécriture :
les joueurs étaient identifiés par leur nom, la collecte consommait vingt fois trop de requêtes, le brut
n'était conservé nulle part en dehors d'un volume Docker, et aucun modèle n'était comparé à une
référence naïve. Surtout, l'abonnement payant à l'API expirait trois semaines plus tard : il fallait
d'abord sauvegarder ce qui ne pourrait plus être récupéré gratuitement.

\section{Objectifs}

Le besoin est double. Le produit est une application locale qui liste les matchs d'une date, dit pour
chacun si les données nécessaires sont disponibles, et affiche alors la loi du nombre total de buts :
probabilités de 0, 1, 2\dots{} buts, espérance, intervalle et probabilité de plus de 2,5 buts, avec les
variables utilisées, la version du modèle et la date des données. L'apprentissage est l'autre moitié :
Git, structuration d'un projet Python, bases de données, statistiques inférentielles, modélisation,
tests, API, interface et documentation. Chaque technologie doit répondre à un besoin concret.

\section{Problématique}

Prédire un nombre de buts est un problème très bruité. Si le total \(T\) suit une loi de Poisson de
paramètre \(\lambda \approx 2{,}8\), la variance irréductible vaut 2,8 alors que la variance de
\(\lambda\) d'un match à l'autre n'est que de l'ordre de 0,3 à 0,5 : même un modèle parfait
n'expliquerait qu'environ 10 à 15~\% de la variance. Juger le projet sur l'erreur d'une prédiction
ponctuelle serait donc trompeur. La question devient : \emph{peut-on prédire une loi de probabilité
mieux calibrée et plus informative qu'une référence simple, et mesurer honnêtement ce qu'apporte chaque
information ?} \og Honnêtement \fg{} veut dire ici : variables disponibles avant le match seulement,
protocole fixé avant les mesures, intervalles de confiance sur chaque comparaison, et un ensemble de
test tenu à l'écart jusqu'à la fin.

\section{Contraintes}

\begin{itemize}
  \item \textbf{Temps} : environ 10 heures par semaine, une personne.
  \item \textbf{Argent} : aucun service payant, hormis l'abonnement API-FOOTBALL déjà payé, terminé le
        22 octobre 2026 (ADR-0005).
  \item \textbf{Matériel} : un portable Windows de 3 à 4 ans ; tout tourne en local.
  \item \textbf{Légalité} : seulement des sources dont les conditions le permettent. Transfermarkt est
        écarté par ses conditions d'utilisation, Understat par son \code{robots.txt} (ADR-0023).
\end{itemize}

\section{Démarche}

Le travail est découpé en six parties (collecte, référentiel, variables, modèles, application,
version intermédiaire), chacune en phase A (développement) puis phase B (exécution sur les données
définitives). Chaque décision est consignée dans une \emph{ADR} (Architecture Decision Record) qui ne se
réécrit pas : 46 ADR à ce jour, listées en annexe~\ref{chap:annexes}. Chaque erreur résolue a son entrée
dans un journal (64 entrées). Le code est écrit avec un assistant de programmation (Claude Code), sous
des règles stables (\code{CLAUDE.md}) : une tâche, une branche, une revue de code, une fusion après
intégration continue verte. Le présent rapport résume ; le détail est dans le dépôt.

\section{Plan du rapport}

Le chapitre~\ref{chap:etat-art} situe les modèles et les métriques dans la littérature. Les
chapitres~\ref{chap:donnees} et~\ref{chap:architecture} décrivent les données et l'architecture. Le
chapitre~\ref{chap:methode} présente les variables, le protocole et les modèles, et le
chapitre~\ref{chap:resultats} leurs résultats. Le chapitre~\ref{chap:application} décrit l'inférence,
l'API et l'interface. Le chapitre~\ref{chap:bilan} fait le bilan : critères de réussite, limites,
difficultés et perspectives.
